\section{Completion of the representation theorem}\label{app:main-proof}

Appendix~\ref{lf:app:pure-trace} proves the one substantial auxiliary input: the pure-trace lemma on the constant-payoff domain.  This appendix combines that result with the affine material scale, calibrates deterministic payoff branches, and completes Theorem~\ref{thm:main}.  We also use the ordinary mixture-space theorem: a complete and transitive mixture order satisfying independence and segment calibration has an affine representation.

The logical order is therefore
\[
\begin{gathered}
\boxed{\text{Axioms \textup{(A1)--(A8)}}}\\[-0.1em]
\begin{array}{ccc}
\swarrow && \searrow\\[-0.4em]
\boxed{\substack{\text{terminal-law quotient and}\\
                  \text{affine material scale}}}
&&
\boxed{\substack{\text{pure-trace lemma (Appendix A) and}\\
                  \text{one common trace scale}}}\\[-0.2em]
\searrow && \swarrow
\end{array}\\[-0.3em]
\boxed{\text{branch calibration and payoff telescope}}\\[-0.1em]
\Downarrow\\[-0.1em]
\boxed{\text{sequentialisation and Theorem 1}}.
\end{gathered}
\]
The left branch identifies the value of terminal payoffs and fixes a common
material scale.  The right branch is proved once on the constant-payoff
domain.  They meet in the deterministic branch calculation, which aligns
material increments with the trace scale and telescopes across all reached
branches.

\subsection{Preliminaries}

We first dispose of three pieces of bookkeeping.  They will be used without
further comment.

Recall that $(q,K)\succeq(p,L)$ denotes the paper's canonical two-block
comparison.  For channels $E,G$ at a common prior $q$, write
$E\succeq_q G$ for $(q,E)\succeq(q,G)$.  A \emph{finite experiment at $q$}
is a finite payoff--record channel paired with that prior; its marked terminal
law is denoted by $\eta_E$.

\begin{lemma}[Blocks, supports, and dummies]\label{lf:lem:bookkeeping}
Suppose Axioms \textup{(A1)--(A5)} hold.
\begin{enumerate}[label=\textup{(\alph*)},leftmargin=*]
\item After identifying alternatives related by invertible action or record
  relabellings, the derived comparison of pairs is complete and transitive.
  If either member of a pair comparison is replaced by an indifferent
  alternative, the comparison is unchanged.
\item If $S=\supp q$, then $(q,K)$ is indifferent to the restriction of
  $(q,K)$ to $S$.
\item If $s\in\D(B)$ and $K^\uparrow(o,r\mid a,b):=K(o,r\mid a)$, then
  \begin{equation}
  (q\otimes s,K^\uparrow)\sim(q,K).
  \label{lf:eq:dummy-ordinal}
  \end{equation}
\end{enumerate}
\end{lemma}

\begin{proof}
Place any finite collection of alternatives in one finite block environment.
Axiom~\textup{(A3)} removes unused blocks, and Axiom~\textup{(A1)} supplies
completeness and transitivity there.  Invertible action and record relabellings
are neutral by applying Axioms~\textup{(A4)} and~\textup{(A5)} in both
directions.  In particular, swapping both tags identifies the two orientations
of a two-block comparison.  This proves part \textup{(a)}; a three-block
environment gives transitivity and a four-block environment gives replacement.

For part \textup{(b)}, report an action in $S$ by its inclusion in $A$.  In the
other direction, report $a\in S$ when the realised action is $a$, and choose an
arbitrary element of $S$ off the support.  Both reports satisfy
\eqref{eq:action-processing}; rows attached to null reported actions may be
completed arbitrarily.  Axiom~\textup{(A5)} gives the two comparisons.

For part \textup{(c)}, first project $(a,b)$ to $a$.  Conversely, given $a$,
draw $b$ independently from $s$ and report $(a,b)$.  Again the two joint-law
identities in \eqref{eq:action-processing} are exact, so Axiom~\textup{(A5)}
applies in both directions.  Part \textup{(a)} completes both arguments.
\end{proof}

Fix for the moment a full-support prior $q\in\D(A)$.  For a joint channel $K$,
write
\[
p_K(o,r):=\sum_a q(a)K(o,r\mid a),
\qquad
\pi^K_{or}(a):=\frac{q(a)K(o,r\mid a)}{p_K(o,r)}
\]
when $p_K(o,r)>0$, and define its \emph{marked terminal law}
\begin{equation}
\eta_K:=\sum_{(o,r):p_K(o,r)>0}
p_K(o,r)\,\delta_{(o,\pi^K_{or})}.
\label{lf:eq:marked-law}
\end{equation}
It records the terminal payoff and the posterior about the realised action.
Write $I_q(E)$ for the mutual information between the action and the complete
visible output of $E$ under $q$.  When the output is a first-stage record $Y$,
we also write this quantity as $I_q(A;Y)$.

\begin{lemma}[Terminal laws and affine fibres]\label{lf:lem:terminal-affine}
Under Axioms \textup{(A1)--(A6)}, the following statements hold.
\begin{enumerate}[label=\textup{(\alph*)},leftmargin=*]
\item If $q$ has full support and $\eta_K=\eta_L$, then
  $(q,K)\sim(q,L)$.
\item At every full-support $q$, the quotient of finite experiments by their
  marked terminal laws has an affine representation $W_q$.
\end{enumerate}
\end{lemma}

\begin{proof}
We give the finite Blackwell argument behind part \textup{(a)}.  For a payoff
$o$ and posterior $\pi$, let
\[
M_K(o,\pi):=
\sum_{\substack{r:\,p_K(o,r)>0\\ \pi^K_{or}=\pi}} p_K(o,r).
\]
Equality of marked laws gives $M_K(o,\pi)=M_L(o,\pi)=:M(o,\pi)$.
For $p_K(o,r)>0$, draw a record $s$ of $L$ according to
\begin{equation}
T(s\mid o,r):=
\1_{\{p_L(o,s)>0,\,\pi^L_{os}=\pi^K_{or}\}}
\frac{p_L(o,s)}{M(o,\pi^K_{or})}.
\label{lf:eq:matching-processor}
\end{equation}
Choose an arbitrary probability row when $p_K(o,r)=0$.  This is a
payoff-preserving record processor.  Bayes' formula gives
\[
q(a)K(o,r\mid a)=p_K(o,r)\pi^K_{or}(a).
\]
For a reached target record $s$, with posterior $\pi=\pi^L_{os}$, it follows
that
\[
\sum_r K(o,r\mid a)T(s\mid o,r)
=\frac{p_L(o,s)}{M(o,\pi)}
  \frac{\pi(a)}{q(a)}M(o,\pi)
=L(o,s\mid a).
\]
If $s$ is unreached, full support makes both sides zero.  Thus postprocessing
$K$ by $T$ gives $L$.  The symmetric construction gives a
processor from $L$ to $K$.  Axiom~\textup{(A4)} and
Lemma~\ref{lf:lem:bookkeeping}\textup{(a)} give indifference.  For a boundary
prior, Lemma~\ref{lf:lem:bookkeeping}\textup{(b)} first deletes the null
actions.

For part \textup{(b)}, take the quotient of actual finite marked experiments by
equality of \eqref{lf:eq:marked-law}.  If $E$ and $G$ are two such experiments,
their public mixture first draws a coin independent of the action, runs the
selected experiment, and retains the coin in the record.  Its marked law is
\begin{equation}
\eta_{tE\oplus(1-t)G}=t\eta_E+(1-t)\eta_G.
\label{lf:eq:public-mixture}
\end{equation}
Axiom~\textup{(A6)} gives, for $0<t<1$ and any third experiment $L$,
\begin{equation}
E\succeq_q L
\quad\Longleftrightarrow\quad
tE\oplus(1-t)G\succeq_q tL\oplus(1-t)G.
\label{lf:eq:mixture-independence}
\end{equation}
For the reverse implication, if $E\not\succeq_q L$, completeness gives
$L\succ_q E$; the strict part of Axiom~\textup{(A6)} makes the reverse public
mixture strict.  If record
alphabets differ, put them in a common disjoint union before invoking the
axiom; part \textup{(a)} identifies the result with
\eqref{lf:eq:public-mixture}.

It remains only to check the Archimedean step.  A closed segment between two
experiments, and any fixed comparator, can be realised on one fixed union of
their finite record alphabets.  Its channel entries vary continuously with
$t\in[0,1]$.
Axiom~\textup{(A2)} therefore makes both contour sets along the segment closed.
Completeness says that they cover $[0,1]$; when a target lies between the two
endpoints, the sets contain opposite endpoints and connectedness makes them
intersect.  Thus every such target is indifferent to a point of the segment.
The mixture-space theorem now supplies an affine representative $W_q$.  This
quotient is closed under every mixture and embedding used below.
\end{proof}

\subsection{The two common scales}

\begin{lemma}[Material utility and a common affine scale]\label{lf:lem:material-scale}
Under Axioms \textup{(A1)--(A7)}, there are outcomes $o^+,o^-$ and a
nonconstant function $u:O\to\R$ such that
\[
u(o^+)=1,
\qquad
u(o^-)=0.
\]
The payoff-lottery order is represented by $\sum_o\ell(o)u(o)$.  The
representatives in Lemma~\ref{lf:lem:terminal-affine} may be normalised so that
\begin{equation}
W_q(\ell)=\sum_o\ell(o)u(o).
\label{lf:eq:payoff-normalisation}
\end{equation}
With this normalisation, adjoining a full-support independent dummy to a
full-support prior preserves value, and for full-support $q,q'$,
\begin{equation}
(q,E)\succeq(q',G)
\quad\Longleftrightarrow\quad
W_q(E)\ge W_{q'}(G).
\label{lf:eq:common-scale}
\end{equation}
\end{lemma}

\begin{proof}
On a singleton action set, terminal-law sufficiency identifies an experiment
with its payoff lottery.  Public-mixture independence and continuity, proved as
in Lemma~\ref{lf:lem:terminal-affine}, give an affine utility on $\D(O)$.
Axiom~\textup{(A7)} supplies $o^+\succ o^-$.  Normalise their values to one
and zero and let $u(o)$ be the value of the degenerate lottery at $o$.
Affinity on the finite simplex gives expected utility.

An action report may ignore its input and draw an arbitrary target prior.
Axiom~\textup{(A5)}, applied in both directions, therefore makes this payoff
lottery order independent of the prior and action alphabet.  Normalise each
$W_q$ at the same two deterministic payoffs.  The payoff-lottery embedding is
affine and order-reflecting, so positive-affine uniqueness gives
\eqref{lf:eq:payoff-normalisation}.

Dummy lifting between full-support fibres is likewise an affine order
embedding.  Positive-affine uniqueness initially permits
\[
W_{q\otimes s}(E^\uparrow)=cW_q(E)+d,
\qquad c>0.
\]
The two payoff anchors have values zero and one on both sides; hence $c=1$ and
$d=0$.  To compare $E$ at $q$ with $G$ at $q'$, lift each to the common product
prior $q\otimes q'$, using the other prior as the dummy.  Equation
\eqref{lf:eq:dummy-ordinal}, the replacement part of
Lemma~\ref{lf:lem:bookkeeping}, and the common product representative give
\eqref{lf:eq:common-scale}.
\end{proof}

We now use Proposition~\ref{pt:prop:main}.  Apply it to the
constant-payoff lift
\[
K_P(o,r\mid a)=\1_{\{o=o^-\}}P(r\mid a).
\]
This lift commutes with block formation, record processing, action reporting,
limits on each finite simplex, and compounding.  Thus Axioms
\textup{(A1)} and \textup{(A8)} give the pure weak order and its positive
orientation; Axiom \textup{(A2)} gives continuity; Axiom \textup{(A3)} gives
block coherence; Axioms \textup{(A4)} and \textup{(A5)} give the two data
processing conditions, including every null-row completion; and Axiom
\textup{(A6)} gives branchwise continuation.  These are precisely the
hypotheses of Proposition~\ref{pt:prop:main}.  Its conclusion is that every
within-channel and block-supported comparison between such lifts is ordered by
$I_{qP}(A;R)$.

\begin{lemma}[The trace scale]\label{lf:lem:trace-scale}
Under Axioms \textup{(A1)--(A8)}, there is one number $\lambda>0$, independent
of the prior and all finite alphabets, such that every full-support,
nontrivial action fibre and every experiment $E$ that pays $o^-$ surely satisfy
\begin{equation}
W_q(E)=\lambda I_q(E).
\label{lf:eq:sure-low}
\end{equation}
\end{lemma}

\begin{proof}
Fix a nontrivial full-support $q$.  On the constant-$o^-$ face, $W_q$ is
affine by Lemma~\ref{lf:lem:terminal-affine}.  For any experiment $E$, the
finite entropy identity
\begin{equation}
I_q(E)=\Sh(q)-\sum_{(o,\pi)\in\supp\eta_E}
\eta_E(o,\pi)\Sh(\pi).
\label{lf:eq:posterior-entropy}
\end{equation}
shows that mutual information depends only on the marked terminal law and is
affine under public mixtures.  Proposition~\ref{pt:prop:main} says that $W_q$ and
$I_q$ represent the same nonconstant order on the constant-low face.
Positive-affine uniqueness therefore gives
\[
W_q(K_P)=\lambda_q I_{qP}+b_q,
\qquad \lambda_q>0.
\]
The uninformative experiment has mutual information zero and, by
\eqref{lf:eq:payoff-normalisation}, value $u(o^-)=0$.  Thus $b_q=0$.

The coefficient does not depend on $q$.  Let $q$ and $q'$ be nontrivial
full-support priors.  Lift full revelation at $q$ to $q\otimes q'$ by adding an
independent dummy coordinate.  Normalised value and mutual information are both
unchanged.  Since the common information value is $\Sh(q)>0$, cancellation
gives $\lambda_{q\otimes q'}=\lambda_q$.  Lifting full revelation from the
other coordinate gives $\lambda_{q\otimes q'}=\lambda_{q'}$.  Hence all such
coefficients agree.  Equivalently, fix the uniform prior on a two-point dummy
alphabet and call its coefficient $\lambda$.  Axiom~\textup{(A8)} fixes the
strict orientation, so $\lambda>0$.

Finally, any experiment that pays $o^-$ surely is a pure record experiment
after the constant payoff coordinate is deleted; restoring that coordinate
changes neither its marked terminal law nor its mutual information.  This
proves \eqref{lf:eq:sure-low}.
\end{proof}

\subsection{The branch calculation}

Let $q$ be full support on a nontrivial $A$ and let
$P:A\to\D(Y)$ be a first-stage record channel.  Put
\[
m_y:=\sum_a q(a)P(y\mid a).
\]
For a deterministic payoff profile $g:Y\to O$, let $C(P,g)$ denote the
compound that first runs $P$ and then pays $g(y)$ with no further record.
Write $g_y^o$ for the profile that pays $o$ at $y$ and $o^-$ elsewhere.

\begin{lemma}[The reached-branch payoff increment]\label{lf:lem:branch-increment}
For every $y\in Y$ and $o\in O$,
\begin{equation}
W_q\bigl(C(P,g_y^o)\bigr)
-W_q\bigl(C(P,o^-)\bigr)=m_y u(o).
\label{lf:eq:branch-increment}
\end{equation}
\end{lemma}

\begin{proof}
If $m_y=0$, changing the payoff on branch $y$ does not change the marked
terminal law, so both sides are zero.  Suppose $m_y>0$, and put
\[
r_y(a):=\frac{q(a)P(y\mid a)}{m_y},
\qquad
S_y:=\supp r_y.
\]
Let $\bar r_y$ be the full-support restriction of $r_y$ to $S_y$.  For a
marked continuation $E$ on $S_y$, extend $E$ to $A$ through a support
projection, insert it at branch $y$, and put the uninformative payoff $o^-$ at
every other branch.  Denote the resulting experiment by $J_y(E)$.  If
$a\notin S_y$, then $r_y(a)=0$; since $q(a)>0$ and $m_y>0$, the definition of
$r_y$ gives $P(y\mid a)=0$.  Hence the continuation rows chosen outside $S_y$
are never used on branch $y$.  Support restriction identifies their local
versions as well, so the extension is independent of the chosen projection.

Support restriction and Axiom~\textup{(A6)}, including its strict part, give
\[
E\succeq_{\bar r_y}G
\quad\Longleftrightarrow\quad
J_y(E)\succeq_q J_y(G).
\]
For the reverse implication, failure of the local weak comparison gives a
strict reverse comparison by completeness, and strict branchwise consistency
gives a strict reverse compound comparison.  Different continuation record
alphabets are first put in a common disjoint union.  Moreover, $J_y$ commutes
with public mixtures at the level of marked terminal laws.  Thus affine
uniqueness gives constants $c_y>0,d_y$ such that
\begin{equation}
W_q(J_y(E))=c_y W_{\bar r_y}(E)+d_y.
\label{lf:eq:insertion-affine}
\end{equation}

Suppose first that $S_y$ has at least two elements.  At payoff $o^-$, compare
full revelation, $E^{\mathrm{id}}$, with an uninformative record, $E^0$.
Lemma~\ref{lf:lem:trace-scale} gives the local difference
\[
W_{\bar r_y}(E^{\mathrm{id}})-W_{\bar r_y}(E^0)
=\lambda\Sh(\bar r_y).
\]
The finite mutual-information chain rule for insertion is
\begin{equation}
I_q(J_y(E))=I_q(A;Y)+m_y I_{\bar r_y}(E).
\label{lf:eq:insertion-chain}
\end{equation}
Support extension does not change the last term.  Both inserted experiments
pay $o^-$ surely, so Lemma~\ref{lf:lem:trace-scale} and
\eqref{lf:eq:insertion-chain} make their aggregate value difference
$\lambda m_y\Sh(\bar r_y)$.  Equation
\eqref{lf:eq:insertion-affine} makes the same difference
$c_y\lambda\Sh(\bar r_y)$.  Since $\lambda>0$ and a nontrivial full-support
distribution has positive entropy, $c_y=m_y$.

Locally, deterministic payoffs $o$ and $o^-$ have values $u(o)$ and zero by
\eqref{lf:eq:payoff-normalisation}.  Subtracting their two instances of
\eqref{lf:eq:insertion-affine}, and identifying the resulting insertions with
the two deterministic compounds in the statement, proves
\eqref{lf:eq:branch-increment}.

If $S_y$ is a singleton, append an independent full-support two-point dummy
action to the first-stage action.  The reached posterior becomes the product of
$r_y$ and the dummy prior and hence has nontrivial support.  The branch mass is
still $m_y$, and Lemma~\ref{lf:lem:material-scale} says that both compound
values are unchanged.  Apply the preceding argument in the product problem and
remove the dummy.
\end{proof}

\begin{lemma}[Deterministic branch assembly]\label{lf:lem:deterministic-assembly}
For every deterministic payoff profile $g:Y\to O$,
\begin{equation}
W_q(C(P,g))
=\lambda I_q(A;Y)+\sum_y m_y u(g(y)).
\label{lf:eq:deterministic-assembly}
\end{equation}
\end{lemma}

\begin{proof}
First observe that the increment from changing one branch is independent of
the payoffs on the other branches.  Indeed, for any two backgrounds $g,h$, the
two crossed public half-mixtures
\[
\tfrac12 C(P,g[y\mapsto o])
\oplus\tfrac12 C(P,h[y\mapsto o^-])
\]
and
\[
\tfrac12 C(P,g[y\mapsto o^-])
\oplus\tfrac12 C(P,h[y\mapsto o])
\]
have the same marked terminal law.  Away from $y$ they contain the same two
payoff-posterior atoms; at $y$ they contain $o$ and $o^-$ with the same mass
$m_y/2$.  Affinity of $W_q$ and cancellation of the two halves show that the
two background-specific increments are equal.  Taking
$h$ to be the all-low profile and applying
Lemma~\ref{lf:lem:branch-increment}, then changing the finitely many branches
one at a time, gives
\begin{equation}
W_q(C(P,g))-W_q(C(P,o^-))
=\sum_y m_y u(g(y)).
\label{lf:eq:telescope}
\end{equation}
The all-low compound pays $o^-$ surely and has exactly the information in the
first-stage record.  Lemma~\ref{lf:lem:trace-scale} therefore gives
$W_q(C(P,o^-))=\lambda I_q(A;Y)$.  Substitute this into
\eqref{lf:eq:telescope}.
\end{proof}

\subsection{Proof of the theorem}

\begin{proof}[Proof of Theorem~\ref{thm:main}]
Assume first Axioms \textup{(A1)--(A8)}.  Let $q$ have full support on a
nontrivial action set, and let $K:A\to\D(O\times R)$ be arbitrary.  Put
$Y=O\times R$ and take as first stage
\[
P_K(y\mid a):=K(y\mid a),
\qquad y=(o,r).
\]
After $y=(o,r)$, pay $o$ for sure and reveal nothing further.  Call the
resulting compound $K^{\mathrm{seq}}$.  From $K$, the record processor sends
$(o,r)$ to the sequential record $((o,r),\ast)$.  In the reverse direction it
sends $((o,r),\ast)$ to $r$; inconsistent payoff-record pairs are unreached and
may be completed arbitrarily.  Both processors preserve the actual payoff
coordinate.  Hence Axiom~\textup{(A4)} gives
$(q,K)\sim(q,K^{\mathrm{seq}})$.

\begin{samepage}
Apply Lemma~\ref{lf:lem:deterministic-assembly} with $g(o,r)=o$.  The two
identities needed are
\[
I_q(A;Y)=I_{qK}(A;O,R)
\]
and
\[
\sum_{o,r}p_{qK}(o,r)u(o)=\E_{qK}[u(O)],
\]
respectively.  Together with sequentialisation, they give
\begin{equation}
W_q(K)=\E_{qK}[u(O)]+\lambda I_{qK}(A;O,R).
\label{lf:eq:full-support-value}
\end{equation}
\end{samepage}

If $A$ is a singleton, mutual information is zero and the marked terminal law
is that of the induced payoff lottery.  Equation
\eqref{lf:eq:payoff-normalisation} gives
\eqref{lf:eq:full-support-value} directly.  If $q$ is on the boundary, put
$S=\supp q$ and write $(q^S,K^S)$ for the restriction.  Then
\begin{equation}
(q,K)\sim(q^S,K^S),
\qquad
V(q,K)=V(q^S,K^S).
\label{lf:eq:support-value}
\end{equation}
The first identity is Lemma~\ref{lf:lem:bookkeeping}\textup{(b)}; the second
follows directly from the finite sums defining expected utility and mutual
information.  Thus no representative $W_q$ is assigned to a boundary fibre.

Restrict two arbitrary alternatives to their supports, apply
\eqref{lf:eq:common-scale} and \eqref{lf:eq:full-support-value} there, and use
\eqref{lf:eq:support-value} to transport the comparison back.  This gives
\begin{equation}
(q,K)\succeq(p,L)
\quad\Longleftrightarrow\quad
V(q,K)\ge V(p,L),
\label{lf:eq:pair-value}
\end{equation}
with the same $u$ and $\lambda$.  Axiom~\textup{(A3)}, first by duplication and
then by deletion of irrelevant blocks, turns \eqref{lf:eq:pair-value} into the
within-channel representation and the finite-block moreover clause.  The
construction made $u$ nonconstant and $\lambda>0$.  This proves the forward
implication.

Conversely, suppose that a nonconstant $u$ and a number $\lambda>0$ represent
the family by $V$.  Axiom~\textup{(A1)} follows from the order on $\R$.
Expected utility and mutual information are jointly continuous in $(K,q)$ on
each finite simplex; for the latter, use the entropy formula and $0\log0=0$.
Hence the set of triples $(K,q,p)$ for which $V(q,K)\ge V(p,K)$ is closed,
which is Axiom~\textup{(A2)}.  A block-supported prior has a constant block
label, which changes neither term of $V$; this proves Axiom~\textup{(A3)} and
reduces every pair comparison below to the corresponding two values.

A payoff-preserving record processor leaves the payoff marginal unchanged and
weakly lowers $I(A;O,R)$ by data processing.  This proves
Axiom~\textup{(A4)}.  Under an action report, equation
\eqref{eq:action-processing} again leaves the payoff-record marginal unchanged,
and the joint law forms the Markov chain
\[
(O,R)-A-B.
\]
Thus $I(B;O,R)\le I(A;O,R)$, proving Axiom~\textup{(A5)}.  This argument uses
the undivided joint-law equation and so also covers every completion on a null
reported action.

For a compound channel, the two finite chain rules are
\begin{align*}
\E[u(O)]
  &=\sum_{y:m(y)>0} m(y)\E_{q_yK^y}[u(O)],\\
I(A;Y,O,R)
  &=I(A;Y)+\sum_{y:m(y)>0} m(y)I_{q_yK^y}(A;O,R).
\end{align*}
The first-stage term is common to two continuation profiles.  Their aggregate
value difference is therefore
\[
\sum_{y:m(y)>0} m(y)\bigl[V(q_y,K^y)-V(q_y,L^y)\bigr].
\]
Branchwise weak improvement makes this sum nonnegative, and a strict
improvement on a reached branch makes it positive.  This is precisely
Axiom~\textup{(A6)}.

Because $u$ is nonconstant, two deterministic payoffs are strictly ranked; the
trace term is zero on a singleton action set.  This gives
Axiom~\textup{(A7)}.  At a constant payoff, full revelation has
value $u(o)+\lambda\Sh(q)$ and an uninformative record has value $u(o)$.
Full support on a nontrivial finite action set implies $\Sh(q)>0$; hence
Axiom~\textup{(A8)}.  The same block calculation used above gives the moreover
clause with the original witnesses $u$ and $\lambda$.
\end{proof}
